\begingroup
\setlength{\abovedisplayskip}{7pt}
\setlength{\belowdisplayskip}{7pt}
\setlength{\abovedisplayshortskip}{3pt}
\setlength{\belowdisplayshortskip}{4pt}

% Printed 122; principal PDF 129; Google PDF 133.
\par\clearpage
\noindent\makebox[0pt][l]{\fontsize{13}{16}\selectfont 122}\hfill{\fontsize{8}{10}\selectfont\textbf{BEMERKUNGEN.\quad STEREOGRAFISCHE PROJEKTION DER KUGEL AUF DIE EBENE.}}\hfill\null\par\medskip

{\small
\noindent \(T\) den Azimut von \(\Delta\) bezeichnet. Ist nun \(\zeta_1\) der Winkel zwischen
\(\Delta_1\) und dem Großkreisbogen \(\delta\), der \((\Delta_1,T_1)\) und
\((\Delta_2,T_2)\) verbindet, und \(z_1\) der Winkel zwischen \(\rho_1\) und \(d\),
so ist \(-(\zeta_1-z_1)\) die Reduktion des sphärischen auf den ebenen Azimut.
Beim Übergang von der positiven \(x\)-Achse zur positiven \(y\)-Achse wächst dabei
der Azimut. Setzt man nun zum Beispiel in die Gleichung
\[
\sin(\zeta_1-z_1)=\sin z_1(\sin\zeta_1\operatorname{cotang}z_1-\cos\zeta_1):
\]
\[
\begin{alignedat}{2}
\sin z_1&=\frac1{2r\operatorname{tang}\dfrac{\Delta_1}2\,.\,d}(x_1y_2-x_2y_1),
&\qquad\operatorname{cotang}z_1&=\frac{\rho_1-\rho_2\cos(T_2-T_1)}{\rho_2\sin(T_2-T_1)}\\
\sin\zeta_1&=\frac{\sin(T_2-T_1)\sin\Delta_2}{\sin\delta},
&\cos\zeta_1&=\frac{\cos\Delta_2-\cos\delta\cos\Delta_1}{\sin\delta\sin T_1},
\end{alignedat}
\]
\noindent so erhält man mit
\(\sin\Delta_1\sin\Delta_2\cos(T_2-T_1)=\cos\delta-\cos\Delta_1\cos\Delta_2\):
\[
-\sin(\zeta_1-z_1)=\frac1{2rd}\operatorname{tang}\frac{\delta}2\,.\,(x_2y_1-x_1y_2).
\]
Die Notiz [2] hat \textsc{Gauss} auf das letzte Blatt seines Exemplars des
»Lehrbegriffs der gesammten Mathematik. Aufgesetzt von \textsc{Wencesl.\ Joh.\ Gust.\
Karsten} etc. Der siebente Theil: Die Optik und Perspectiv. Greifswald, 1775«
eingetragen. \textsc{Gauss} hat sie mit der Überschrift »ad pag.\ 718« versehen.
\par\noindent [3] findet sich auf dem letzten Blatt des \textsc{Gauss}schen Exemplars
der »Beyträge zum Gebrauche der Mathematik etc. von \textsc{J.\ H.\ Lambert.}
Dritter Theil. Berlin, 1772.«
\par
\noindent\hfill\textsc{Krüger, Börsch.}\hspace*{.06\linewidth}\par}
\vspace{2\baselineskip}
\begin{center}\rule{.24\linewidth}{.4pt}\end{center}

% Printed 123; principal PDF 130; Google PDF 134.
\par\clearpage

\thispagestyle{empty}
\vspace*{4\baselineskip}
\begin{center}
{\large [III.]}\par\bigskip
{\Large [ÜBERTRAGUNG DER KUGEL AUF DIE EBENE\\
DURCH MERCATORS PROJEKTION.]}\par\bigskip
\rule{.09\linewidth}{.4pt}\par\bigskip
[1.]
\end{center}
Erstreckt sich die Vermessung hauptsächlich entlang des Meridians und nur
wenig nach Osten und Westen, so ist \textsc{Mercators} Projektion zweckmäßiger,
wenn der Meridian als Hauptkreis betrachtet wird.

Die Hauptaufgaben sind hier die folgenden.

[Es sei
\begin{quote}
\setlength{\tabcolsep}{0pt}
\begin{tabular}{@{}lp{.82\linewidth}@{}}
\(r\)\enspace.\,.\,.\,.\,.\,. & Radius der Kugel\\
\(P\)\enspace.\,.\,.\,.\,. & geografische Breite des Anfangspunkts auf dem Hauptmeridian\\
\(\psi,\lambda\)\,.\,.\,.\,. & geografische Breite und Länge eines beliebigen Punkts der
Kugeloberfläche, dessen rechtwinklige sphärische Koordinaten \(X,Y\) sind\\
\(x,y\)\enspace.\,.\,. & rechtwinklige Koordinaten des entsprechenden Punkts in der Ebene.
\(x\) ist nach Süden positiv, \(y\) wächst mit dem Längenunterschied \(\lambda\).
\quad \(X=\frac xr,\ \operatorname{tang}(45^\circ+\tfrac12Y)=e^{\frac yr}\).]
\end{tabular}
\end{quote}
I.\quad Korrektur der berechneten Linienrichtung:
\[
\frac{(x'-x)(y'+2y)}{6rr}.
\]
\noindent\hfill \(16^{*}\)\hspace*{.14\linewidth}\par

% Printed 124; principal PDF 131; Google PDF 135.
\par\clearpage
\noindent\makebox[0pt][l]{\fontsize{13}{16}\selectfont 124}\hfill{\small\textbf{NACHLASS.}}\hfill\null\par\medskip

II.\quad Wahre Entfernung:
\[
\surd\left\{\left((x'-x)^2+(y'-y)^2\right)
 \left(1-\tfrac16\frac{yy+yy'+y'y'}{rr}\right)\right\}.
\]
III.\quad Geografische Lage:
\[
\begin{aligned}
\operatorname{tang}\lambda&=\frac{y\left(1+\dfrac{yy}{6rr}\right)}{r\cos\left(P-\dfrac xr\right)}\\
\operatorname{tang}\text{Meridiankonv.}
 &=\frac yr\left(1-\frac{yy}{3rr}\right)\operatorname{tang}\left(P-\frac xr\right)\\
\sin\text{Meridiankonv.}&=\sin\lambda\sin\left(P-\frac xr\right)\\
\operatorname{tang}\tfrac12\left(P-\frac xr-\psi\right)
 &=\tfrac12\frac yr\left(1-\frac{yy}{12rr}\right)\operatorname{tang}\tfrac12
 \text{Meridiankonv.}
\end{aligned}
\]
IV.\quad Die Umkehraufgabe:
\[
\begin{aligned}
y&=r\sin\lambda\cos\psi\,.\,\left(1+\frac{yy}{3rr}\right)\\
\sin\text{Meridiankonv.}&=\frac yr\left(1+\frac{yy}{6rr}\right)\operatorname{tang}\psi\\
 &=\sin\lambda\sin\psi\,.\,\left(1+\frac{yy}{2rr}\right)\\
\operatorname{tang}\text{Meridiankonv.}&=\operatorname{tang}\lambda\sin\psi.
\end{aligned}
\]

\begin{center}[2.]\end{center}
Konvergenz der Meridiane \(=c\).
\[
\begin{aligned}
\operatorname{tang}Y&=\frac yr+\tfrac16\left(\frac yr\right)^3+
 \tfrac1{120}\left(\frac yr\right)^5\ldots\\
\sec Y&=1+\tfrac12\frac{yy}{rr}+\tfrac1{24}\left(\frac yr\right)^4\ldots\\
\sin Y&=\frac yr-\tfrac13\left(\frac yr\right)^3+
 \left[\tfrac2{15}\left(\frac yr\right)^5\right]\ldots\\
\operatorname{tang}\lambda&=
 \frac{\dfrac yr+\tfrac16\left(\dfrac yr\right)^3+\tfrac1{120}\left(\dfrac yr\right)^5\ldots}
 {\cos\left(P-\dfrac xr\right)}\\
\operatorname{tang}c&=\operatorname{tang}\left(P-\frac xr\right)\,.\,
 \left(\frac yr-\tfrac13\left(\frac yr\right)^3+\tfrac2{15}\left(\frac yr\right)^5\ldots\right)\\
\sin c&=\sin\lambda\sin\left(P-\frac xr\right).
\end{aligned}
\]
\noindent [Setzt man]
\[
P-\frac xr-\psi=R,
\]

% Printed 125; principal PDF 132; Google PDF 136.
\par\clearpage
\noindent\null\hfill{\fontsize{8}{10}\selectfont\textbf{ÜBERTRAGUNG DER KUGEL AUF DIE EBENE DURCH MERCATORS PROJEKTION.}}\hfill\makebox[0pt][r]{\fontsize{13}{16}\selectfont 125}\par\medskip

\noindent [so ergibt sich]
\[
\begin{aligned}
\sin\tfrac12R
 &=\frac{\sin\tfrac12 c}{\cos\tfrac12\lambda}\,.\,\tfrac12\frac yr
 \left(1-\tfrac5{24}\frac{yy}{rr}+\tfrac{121}{1920}\left(\frac yr\right)^4\ldots\right)\\
\cos\tfrac12R
 &=\frac{\cos\tfrac12 c}{\cos\tfrac12\lambda}
 \left(1-\tfrac18\frac{yy}{rr}+\tfrac{17}{384}\left(\frac yr\right)^4\ldots\right)\\
\operatorname{tang}\tfrac12R
 &=\operatorname{tang}\tfrac12c\,.\,\tfrac12\frac yr
 \left(1-\tfrac1{12}\frac{yy}{rr}+\tfrac1{120}\left(\frac yr\right)^4\ldots\right)\\
 &=\operatorname{tang}\tfrac12Y^2\operatorname{tang}\left(P-\frac xr-\tfrac12R\right)\\
\sin R&=\frac{\sin c}{\cos\tfrac12\lambda^2}\,.\,\tfrac12\frac yr
 \left(1-\tfrac13\frac{yy}{rr}+\tfrac2{15}\left(\frac yr\right)^4\ldots\right)\\
 &=\tfrac12\frac{\sin c\operatorname{tang}c}{\cos\tfrac12\lambda^2}
 \operatorname{cotang}\left(P-\frac xr\right)\\
 &=\operatorname{tang}\tfrac12\lambda\operatorname{tang}c\cos\left(P-\frac xr\right)\\
 &=\frac{\sin\tfrac12\lambda^2\sin2\left(P-\dfrac xr\right)}{\cos c}\\
 &=\tfrac12\left(\frac{yy}{rr}\ldots\right)\operatorname{tang}\left(P-\frac xr\right)
\end{aligned}
\]
\[
R=\tfrac12\frac{yy}{rr}\operatorname{tang}\left(P-\frac xr\right)
 -\left(\frac yr\right)^4\left(\tfrac18\operatorname{tang}\left(P-\frac xr\right)^2+\tfrac5{24}\right)\ldots;
\]
\noindent [für \(r=1\) gilt]
\[
\frac{\partial R}{\partial y}=
\frac{\sin(P-x)\,.\,(y-\tfrac13y^3+\tfrac2{15}y^5\ldots)}
{\surd(\cos(P-x)^2+yy+\tfrac13y^4+\tfrac2{45}y^6\ldots)}.
\]

\begin{center}
[3.]\par\medskip
Hilfstafel für \(\log\sec Y\).
\end{center}
Für \(r=1\) ergibt sich [der Logarithmus des Maßstabsfaktors]
\[
\log\sec Y=\tfrac12 yy-\tfrac1{12}y^4+\tfrac1{45}y^6-\tfrac{17}{2520}y^8\ldots.
\]
Beim Wert von \(r\) für Göttingen [mit \(P=51^\circ31^{\prime}48^{\prime\prime}{,}70\)
und der Abplattung \(\frac1{302{,}68}\) ergibt sich aus
\(r=\frac a{\surd(1-ee\sin P^2)}\): \(\log r=6{,}805\,4777\)] lauten
die [Briggs'schen] Logarithmen der Koeffizienten auf 7 Dezimalstellen:
\[
2{,}725\,7989-10\qquad 8{,}336\,6923-30\qquad
4{,}151\,7056-40\qquad 0{,}023\,0111-50.
\]
[Damit wurde die Tabelle auf der nächsten Seite berechnet.]

% Printed 126; principal PDF 133; Google PDF 137.
\par\clearpage
\noindent\makebox[0pt][l]{\fontsize{13}{16}\selectfont 126}\hfill{\small\textbf{NACHLASS.}}\hfill\null\par\medskip

\vspace{2.4\baselineskip}
{\small
\setlength{\tabcolsep}{6pt}
\renewcommand{\arraystretch}{1.12}
\noindent\makebox[\linewidth][c]{%
\begin{tabular}{|r|r|r||r|r|r||r|r|r|}
\hline
\multicolumn{1}{|c|}{\llap{[\hspace{1.15em}}\(y\)} & \raisebox{0pt}[\dimexpr\height+2pt\relax]{\shortstack{\(\log\sec Y\)\\in Einh.\\d. 7. Dez.}} &
\multicolumn{1}{c||}{diff.} &
\multicolumn{1}{c|}{\(y\)} & \raisebox{0pt}[\dimexpr\height+2pt\relax]{\shortstack{\(\log\sec Y\)\\in Einh.\\d. 7. Dez.}} &
\multicolumn{1}{c||}{diff.} &
\multicolumn{1}{c|}{\(y\)} & \raisebox{0pt}[\dimexpr\height+2pt\relax]{\shortstack{\(\log\sec Y\)\\in Einh.\\d. 7. Dez.}} &
\multicolumn{1}{c|}{diff.\rlap{\hspace{.85em}]}} \\
\hline
\rule{0pt}{.5em}&&&&&&&&\\[-.6ex]
\multicolumn{1}{|c|}{0} & \multicolumn{1}{c|}{0} & \raisebox{-6.7pt}[0pt][0pt]{5.32} & 250\,000 & 3323.29 & \raisebox{-6.7pt}[0pt][0pt]{271.10} & 500\,000 & 13283.00 & \raisebox{-6.7pt}[0pt][0pt]{536.06} \\
10\,000 & 5.32 & \raisebox{-6.7pt}[0pt][0pt]{15.95} & 260\,000 & 3594.39 & \raisebox{-6.7pt}[0pt][0pt]{281.73} & 510\,000 & 13819.06 & \raisebox{-6.7pt}[0pt][0pt]{546.64} \\
20\,000 & 21.27 & \raisebox{-6.7pt}[0pt][0pt]{26.60} & 270\,000 & 3876.12 & \raisebox{-6.7pt}[0pt][0pt]{292.34} & 520\,000 & 14365.70 & \raisebox{-6.7pt}[0pt][0pt]{557.20} \\
30\,000 & 47.87 & \raisebox{-6.7pt}[0pt][0pt]{37.23} & 280\,000 & 4168.46 & \raisebox{-6.7pt}[0pt][0pt]{302.96} & 530\,000 & 14922.90 & \raisebox{-6.7pt}[0pt][0pt]{567.77} \\
40\,000 & 85.10 & \raisebox{-6.7pt}[0pt][0pt]{47.86} & 290\,000 & 4471.42 & \raisebox{-6.7pt}[0pt][0pt]{313.58} & 540\,000 & 15490.67 & \raisebox{-6.7pt}[0pt][0pt]{578.33} \\
50\,000 & 132.96 & \raisebox{-6.7pt}[0pt][0pt]{58.51} & 300\,000 & 4785.00 & \raisebox{-6.7pt}[0pt][0pt]{324.19} & 550\,000 & 16069.00 & \raisebox{-6.7pt}[0pt][0pt]{588.88} \\
\cline{1-2}\cline{4-5}\cline{7-8}
60\,000 & 191.47 & \raisebox{-6.7pt}[0pt][0pt]{69.14} & 310\,000 & 5109.19 & \raisebox{-6.7pt}[0pt][0pt]{334.80} & 560\,000 & 16657.88 & \raisebox{-6.7pt}[0pt][0pt]{599.44} \\
70\,000 & 260.61 & \raisebox{-6.7pt}[0pt][0pt]{79.77} & 320\,000 & 5443.99 & \raisebox{-6.7pt}[0pt][0pt]{345.41} & 570\,000 & 17257.32 & \raisebox{-6.7pt}[0pt][0pt]{610.00} \\
80\,000 & 340.38 & \raisebox{-6.7pt}[0pt][0pt]{90.41} & 330\,000 & 5789.40 & \raisebox{-6.7pt}[0pt][0pt]{356.02} & 580\,000 & 17867.32 & \raisebox{-6.7pt}[0pt][0pt]{620.55} \\
90\,000 & 430.79 & \raisebox{-6.7pt}[0pt][0pt]{101.05} & 340\,000 & 6145.42 & \raisebox{-6.7pt}[0pt][0pt]{366.63} & 590\,000 & 18487.87 & \raisebox{-6.7pt}[0pt][0pt]{631.09} \\
100\,000 & 531.84 & \raisebox{-6.7pt}[0pt][0pt]{111.69} & 350\,000 & 6512.05 & \raisebox{-6.7pt}[0pt][0pt]{377.24} & 600\,000 & 19118.96 & \raisebox{-6.7pt}[0pt][0pt]{641.63} \\
\cline{1-2}\cline{4-5}\cline{7-8}
110\,000 & 643.53 & \raisebox{-6.7pt}[0pt][0pt]{122.31} & 360\,000 & 6889.29 & \raisebox{-6.7pt}[0pt][0pt]{387.83} & 610\,000 & 19760.59 & \raisebox{-6.7pt}[0pt][0pt]{652.18} \\
120\,000 & 765.84 & \raisebox{-6.7pt}[0pt][0pt]{132.95} & 370\,000 & 7277.12 & \raisebox{-6.7pt}[0pt][0pt]{398.45} & 620\,000 & 20412.77 & \raisebox{-6.7pt}[0pt][0pt]{662.71} \\
130\,000 & 898.79 & \raisebox{-6.7pt}[0pt][0pt]{143.58} & 380\,000 & 7675.57 & \raisebox{-6.7pt}[0pt][0pt]{409.04} & 630\,000 & 21075.48 & \raisebox{-6.7pt}[0pt][0pt]{673.25} \\
140\,000 & 1042.37 & \raisebox{-6.7pt}[0pt][0pt]{154.21} & 390\,000 & 8084.61 & \raisebox{-6.7pt}[0pt][0pt]{419.63} & 640\,000 & 21748.73 & \raisebox{-6.7pt}[0pt][0pt]{683.78} \\
150\,000 & 1196.58 & \raisebox{-6.7pt}[0pt][0pt]{164.84} & 400\,000 & 8504.24 & \raisebox{-6.7pt}[0pt][0pt]{430.23} & 650\,000 & 22432.51 & \raisebox{-6.7pt}[0pt][0pt]{694.31} \\
\cline{1-2}\cline{4-5}\cline{7-8}
160\,000 & 1361.42 & \raisebox{-6.7pt}[0pt][0pt]{175.48} & 410\,000 & 8934.47 & \raisebox{-6.7pt}[0pt][0pt]{440.83} & 660\,000 & 23126.82 & \raisebox{-6.7pt}[0pt][0pt]{704.83} \\
170\,000 & 1536.90 & \raisebox{-6.7pt}[0pt][0pt]{186.10} & 420\,000 & 9375.30 & \raisebox{-6.7pt}[0pt][0pt]{451.42} & 670\,000 & 23831.65 & \raisebox{-6.7pt}[0pt][0pt]{715.36} \\
180\,000 & 1723.00 & \raisebox{-6.7pt}[0pt][0pt]{196.74} & 430\,000 & 9826.72 & \raisebox{-6.7pt}[0pt][0pt]{462.00} & 680\,000 & 24547.01 & \raisebox{-6.7pt}[0pt][0pt]{725.88} \\
190\,000 & 1919.74 & \raisebox{-6.7pt}[0pt][0pt]{207.36} & 440\,000 & 10288.72 & \raisebox{-6.7pt}[0pt][0pt]{472.59} & 690\,000 & 25272.89 & \raisebox{-6.7pt}[0pt][0pt]{736.38} \\
200\,000 & 2127.10 & \raisebox{-6.7pt}[0pt][0pt]{217.99} & 450\,000 & 10761.31 & \raisebox{-6.7pt}[0pt][0pt]{483.18} & 700\,000 & 26009.27 & \raisebox{-6.7pt}[0pt][0pt]{746.90} \\
\cline{1-2}\cline{4-5}\cline{7-8}
210\,000 & 2345.09\llap{\raisebox{-.85ex}[0pt][0pt]{.}} & \raisebox{-6.7pt}[0pt][0pt]{228.61} & 460\,000 & 11244.49 & \raisebox{-6.7pt}[0pt][0pt]{493.76} & 710\,000 & 26756.17 & \raisebox{-6.7pt}[0pt][0pt]{757.40} \\
220\,000 & 2573.70 & \raisebox{-6.7pt}[0pt][0pt]{239.24} & 470\,000 & 11738.25 & \raisebox{-6.7pt}[0pt][0pt]{504.34} & 720\,000 & 27513.57 & \raisebox{-6.7pt}[0pt][0pt]{767.91} \\
230\,000 & 2812.94 & \raisebox{-6.7pt}[0pt][0pt]{249.86} & 480\,000 & 12242.59 & \raisebox{-6.7pt}[0pt][0pt]{514.92} & 730\,000 & 28281.48 & \raisebox{-6.7pt}[0pt][0pt]{778.41} \\
240\,000 & 3062.80 & \raisebox{-6.7pt}[0pt][0pt]{260.49} & 490\,000 & 12757.51 & \raisebox{-6.7pt}[0pt][0pt]{525.49} & 740\,000 & 29059.89 & \raisebox{-6.7pt}[0pt][0pt]{788.90} \\
250\,000 & 3323.29 &  & 500\,000 & 13283.00 &  & 750\,000 & 29848.79 &  \\
\end{tabular}%
}\par}
% Printed 127; principal PDF 134; Google PDF 138.
\par\clearpage
\noindent\null\hfill{\fontsize{8}{10}\selectfont\textbf{ÜBERTRAGUNG DER KUGEL AUF DIE EBENE DURCH MERCATORS PROJEKTION.}}\hfill\makebox[0pt][r]{\fontsize{13}{16}\selectfont 127}\par\medskip

\begin{center}[4.]\end{center}
[Ist der] Radius [der Kugel] \(=1\), [so gilt]
\[
\begin{aligned}
y&=\log\operatorname{tang}(45^\circ+\tfrac12Y)\\
\operatorname{tang}\tfrac12iy&=i\operatorname{tang}\tfrac12Y\\
\cos\tfrac12iy&=\frac{\cos\tfrac12Y}{\surd\cos Y}=\cos\tfrac12Y\surd\cos iy\\
\sin iy&=i\operatorname{tang}Y\\
\operatorname{tang}iy&=i\sin Y\\
\cos iy&=\frac1{\cos Y}=\text{Maßstabsfaktor.}
\end{aligned}
\]

\begin{center}[5.]\end{center}
[Es sei \(Q\) die Polhöhe an einem beliebigen Punkt der Kugeloberfläche
mit der Ordinate \(Y\) und \(Q^{*}\) die Polhöhe am Fußpunkt von \(Y\).
Ferner sei \(c\) die Meridiankonvergenz. Dann gilt:]
\[
\begin{aligned}
\operatorname{tang}\lambda&=\frac{\sin iy}{i\cos Q^{*}}=\frac{\operatorname{tang}c}{\sin Q}\\
\sin\lambda&=\frac{\operatorname{tang}iy}{i\cos Q}=\frac{\sin c}{\sin Q^{*}}\\
\operatorname{tang}c&=\frac{\operatorname{tang}iy}{i}\operatorname{tang} Q^{*}\\
\sin c&=\sin\lambda\sin Q^{*}\\
\cos c&=\cos\lambda\cos iy\\
\operatorname{tang}\tfrac12(Q^{*}-Q)&=\operatorname{tang}\tfrac12c
 \frac{\operatorname{tang}\tfrac12iy}{i}.
\end{aligned}
\]

\begin{center}[6.]\end{center}
Um aus
\[
\operatorname{tang}iy\,[=i\sin Y]
\]
\noindent \(y\) zu bestimmen, setze man
\[
\sin Y=\operatorname{tang}(Y-\delta);
\]
\noindent dann gilt [näherungsweise]
\[
y=\frac{\sin Y}{\cos(Y-\tfrac1{15}\delta)^{\frac23}}.
\]

% Printed 128; principal PDF 135; Google PDF 139.
\par\clearpage
\noindent\makebox[0pt][l]{\fontsize{13}{16}\selectfont 128}\hfill{\small\textbf{NACHLASS.}}\hfill\null\par\medskip

[Denn setzt man \(\sin Y=s\), so gilt:]
\[
\begin{aligned}
Y&=s+\tfrac16s^3+\tfrac3{40}s^5+\tfrac5{112}s^7\ldots\\
Y-\delta&=s-\tfrac13s^3+\tfrac15s^5-\tfrac17s^7\ldots,
\end{aligned}
\]
\noindent also
\[
\begin{aligned}
Y-\tfrac1{15}\delta&=s+\tfrac2{15}s^3+\tfrac1{12}s^5+\tfrac9{280}s^7\ldots\\
\cos(Y-\tfrac1{15}\delta)&=1-\tfrac12ss-\tfrac{11}{120}s^4-\tfrac{257}{3600}s^6\ldots\\
\cos(Y-\tfrac1{15}\delta)^{-\frac23}&=1+\tfrac13ss+\tfrac15s^4+\tfrac{649}{4050}s^6\ldots,
\end{aligned}
\]
\noindent [und somit
\[
\frac{\sin Y}{\cos(Y-\tfrac1{15}\delta)^{\frac23}}
=\sin Y+\tfrac13\sin Y^3+\tfrac15\sin Y^5+\tfrac{649}{4050}\sin Y^7\ldots.
\]
\noindent Genau gilt
\[
y=\sin Y+\tfrac13\sin Y^3+\tfrac15\sin Y^5+\tfrac17\sin Y^7\ldots\mathclose{]}.
\]
\medskip
Zur Herleitung von \(\log y\) aus \(\log\sin Y\) [ergibt die letzte Gleichung]
\[
\log y=\log\sin Y+\tfrac13\sin Y^2+\tfrac{13}{90}\sin Y^4+
 \tfrac{251}{2835}\sin Y^6\ldots.
\]

\begin{center}
[7.]\par\medskip
\textsc{Mercators} Projektion.
\end{center}
\begin{quote}
\setlength{\tabcolsep}{0pt}
\begin{tabular}{@{}lp{.83\linewidth}@{}}
\(x,x'\)\,.\,.\,.\,.\,. &
\makebox[5.6em][l]{Abszissen}\raisebox{-.5\baselineskip}[0pt][0pt]{\(\left.\vphantom{\begin{matrix}x\\x\end{matrix}}\right\}\)}
\raisebox{-.5\baselineskip}[0pt][0pt]{\enspace[zweier Punkte in der Ebene]}\\
\(y,y'\)\,.\,.\,.\,.\,. & \makebox[5.6em][l]{Ordinaten}\\
\(Y,Y'\)\,.\,.\,.\,. & den Ordinaten entsprechende Bögen [auf der Kugel]\\
\(D\)\,.\,.\,.\,.\,.\,.\,. & Entfernung der beiden Punkte [auf der Kugel]\\
{[}\(d\)\,.\,.\,.\,.\,.\,.\,.\,. & Entfernung der beiden Punkte in der Ebene]\\
\(A,A'\)\,.\,.\,.\,. & Neigungswinkel von \(D\) oder von den Tangenten an seinen beiden
Endpunkten zur Abszissenlinie
\end{tabular}

Radius [der Kugel] \(=1\).
\end{quote}
Bei der Darstellung der Kugeloberfläche in der Ebene nach \textsc{Mercators}
Projektion gelten die folgenden Beziehungen:
\[
\begin{aligned}
\sin\tfrac12(A+A')\sin\tfrac12D&=\cos\tfrac12(x'-x)\sin\tfrac12(Y'-Y)\\
\cos\tfrac12(A+A')\sin\tfrac12D&=\sin\tfrac12(x'-x)\cos\tfrac12(Y'+Y)\\
\sin\tfrac12(A-A')\cos\tfrac12D&=\sin\tfrac12(x'-x)\sin\tfrac12(Y'+Y)\\
\cos\tfrac12(A-A')\cos\tfrac12D&=\cos\tfrac12(x'-x)\cos\tfrac12(Y'-Y)
\end{aligned}
\]

% Printed 129; principal PDF 136; Google PDF 140.
\par\clearpage
\noindent\null\hfill{\fontsize{8}{10}\selectfont\textbf{ÜBERTRAGUNG DER KUGEL AUF DIE EBENE DURCH MERCATORS PROJEKTION.}}\hfill\makebox[0pt][r]{\fontsize{13}{16}\selectfont 129}\par\medskip

\noindent [oder, da
\[
\sin\tfrac12Y=\frac1i\frac{\sin\frac{iy}2}{\surd\cos iy},\qquad
\cos\tfrac12Y=\frac{\cos\frac{iy}2}{\surd\cos iy}\quad\text{gilt,]}
\]
\[
\begin{aligned}
i\sin\tfrac12(A+A')\sin\tfrac12D\,.\,\surd(\cos iy'\cos iy)
 &=\cos\tfrac12(x'-x)\sin\tfrac12 i(y'-y)\\
\cos\tfrac12(A+A')\sin\tfrac12D\,.\,\surd(\cos iy'\cos iy)
 &=\sin\tfrac12(x'-x)\cos\tfrac12 i(y'-y)\\
i\sin\tfrac12(A-A')\cos\tfrac12D\,.\,\surd(\cos iy'\cos iy)
 &=\sin\tfrac12(x'-x)\sin\tfrac12 i(y'+y)\\
\cos\tfrac12(A-A')\cos\tfrac12D\,.\,\surd(\cos iy'\cos iy)
 &=\cos\tfrac12(x'-x)\cos\tfrac12 i(y'+y);
\end{aligned}
\]
\noindent [und daraus]
\[
\begin{aligned}
\operatorname{tang}\tfrac12(A+A')&=
\frac{\operatorname{tang}\tfrac12 i(y'-y)}{i\operatorname{tang}\tfrac12(x'-x)}\\
\operatorname{tang}\tfrac12(A-A')&=\frac1i\operatorname{tang}\tfrac12(x'-x)
 \operatorname{tang}\tfrac12 i(y'+y).
\end{aligned}
\]
\noindent [Da]
\[
\log\operatorname{tang}u=\log u+\tfrac13uu+\tfrac7{90}u^4+\tfrac{62}{2835}u^6\ldots
\]
\noindent [gilt, ergibt sich zunächst
\[
\log\operatorname{tang}\tfrac12(A+A')=\log\operatorname{tang}\alpha
 -\tfrac1{12}dd-\tfrac7{1440}d^4\cos2\alpha
 -\tfrac{31}{725760}(5+3\cos4\alpha)\ldots,
\]
\noindent wobei]
\[
\begin{aligned}
x'-x&=d\cos\alpha\\
y'-y&=d\sin\alpha
\end{aligned}
\]
\noindent [gesetzt wird].

Wenn nun
\[
\log\operatorname{tang}(t+G)=\log\operatorname{tang}t+g
\]
\noindent [gilt], so folgt
\[
\begin{aligned}
\frac{\mathrm dG}{\mathrm dg}&=\tfrac12\sin2(t+G)\\
\frac{\mathrm d\mathrm dG}{\mathrm dg^2}&=\tfrac14\sin4(t+G)\\
\frac{\mathrm d^3G}{\mathrm dg^3}&=\tfrac14\sin6(t+G)-\tfrac14\sin2(t+G)
\end{aligned}
\]
\centerline{usw.,}
\medskip
\noindent [folglich]
\[
\begin{aligned}
G&=\tfrac12g\sin2t+\tfrac18gg\sin4t+\tfrac1{24}g^3(\sin6t-\sin2t)\ldots\\
 &=\tfrac12g\sin2t(1+\tfrac12g\cos2t+\tfrac16gg\cos4t\ldots).
\end{aligned}
\]
\noindent\textsc{ix.}\hfill 17\hspace*{.13\linewidth}\par

% Printed 130; principal PDF 137; Google PDF 141.
\par\clearpage
\noindent\makebox[0pt][l]{\fontsize{13}{16}\selectfont 130}\hfill{\small\textbf{NACHLASS.}}\hfill\null\par\medskip

Daraus ergibt sich [wenn
\[
G=\tfrac12(A+A')-\alpha,\qquad t=\alpha,\qquad
g=-\tfrac1{12}dd-\tfrac7{1440}d^4\cos2\alpha-\ldots
\]
\noindent gesetzt wird:]
\[
\begin{aligned}
\tfrac12(A+A')&=\alpha-\tfrac1{24}dd\sin2\alpha-\tfrac1{2880}d^4\sin4\alpha
 -\tfrac1{181440}d^6\sin6\alpha\ldots\\
 &=\alpha-\tfrac1{12}(x'-x)(y'-y)
 -\tfrac1{720}(x'-x)(y'-y)((x'-x)^2-(y'-y)^2)\ldots.
\end{aligned}
\]
Aus
\[
\operatorname{tang}u=\operatorname{tang}v\frac{\operatorname{tang}iw}{i}
\]
\noindent folgt
\[
u=vw\{1+\tfrac13(vv-ww)+\tfrac2{45}(3vv-ww)(vv-3ww)\ldots\},
\]
\noindent [somit ergibt sich für \(u=\tfrac12(A-A')\), \(v=\tfrac12(x'-x)\),
\(w=\tfrac12(y'+y)\):]
\[
\tfrac12(A-A')=\tfrac14(x'-x)(y'+y)
 +\tfrac1{48}(x'-x)(y'+y)((x'-x)^2-(y'+y)^2)\ldots.
\]
\noindent [Wird]
\[
\begin{aligned}
x'-x&=a\\
y'-y&=b\\
y'+y&=c
\end{aligned}
\]
\noindent [gesetzt, so gilt auch]
\[
\begin{aligned}
\tfrac12(A+A')&=\alpha-\tfrac1{12}ab-\tfrac1{720}ab(aa-bb)\ldots\\
\tfrac12(A-A')&=\tfrac14ac-\tfrac1{48}ac(cc-aa)\ldots\\
A&=\alpha+\frac{2y+y'}6a-\tfrac1{720}a(15c^3-b^3)+\tfrac1{720}a^3(15c-b)\ldots\\
\mathopen{[}A'&=\alpha-\frac{y+2y'}6a+\tfrac1{720}a(15c^3+b^3)
 -\tfrac1{720}a^3(15c+b)\ldots\mathclose{]}.
\end{aligned}
\]

\begin{center}
[8.]\par\medskip
Reduktion der Hilfskugel auf die Ebene.
\end{center}
[Setzt man
\[
e^{\frac12i(x'-x)}=\xi,
\]
\noindent so folgt aus den zuerst aufgeführten Gleichungen in Abschnitt 7 unten auf S.\ 128:]
\[
\sin\tfrac12D\,.\,e^{\frac12i(A+A')}
 =\frac{\xi\xi-1}{2i\xi}\cos\tfrac12(Y'+Y)
 +\frac{i(\xi\xi+1)}{2\xi}\sin\tfrac12(Y'-Y)
\]

% Printed 131; principal PDF 138; Google PDF 142.
\par\clearpage
\noindent\null\hfill{\fontsize{8}{10}\selectfont\textbf{ÜBERTRAGUNG DER KUGEL AUF DIE EBENE DURCH MERCATORS PROJEKTION.}}\hfill\makebox[0pt][r]{\fontsize{13}{16}\selectfont 131}\par\medskip

\noindent [oder, wenn
\[
Y'=90^\circ-p',\qquad Y=90^\circ-p
\]
\noindent gesetzt wird,]
\[
\begin{aligned}
\sin\tfrac12D\,.\,e^{\frac12i(A+A')}
 &=-\tfrac12i\xi\{\sin\tfrac12(p+p')-\sin\tfrac12(p-p')\}\\
 &\quad+\tfrac12i\xi^{-1}\{\sin\tfrac12(p+p')+\sin\tfrac12(p-p')\}\\
 &=-i\xi\cos\tfrac12p\sin\tfrac12p'+i\xi^{-1}\sin\tfrac12p\cos\tfrac12p'
\end{aligned}
\]
\noindent [oder, wegen]
\[
e^y=\operatorname{cotang}\tfrac12p,
\]
\[
\begin{aligned}
\sin\tfrac12D\,.\,e^{\frac12i(A+A')}
 &=\frac{-i\xi}{\surd\{(1+e^{2y'})(1+e^{-2y})\}}
 +\frac{i\xi^{-1}}{\surd\{(1+e^{-2y'})(1+e^{2y})\}}\\
 &=\frac{-i\xi\,.\,e^{\frac12(y-y')}+i\xi^{-1}\,.\,e^{\frac12(y'-y)}}
 {\surd\{(e^y+e^{-y})(e^{y'}+e^{-y'})\}}\\
 &=\frac{-ie^{\frac12(y-y')+\frac12i(x'-x)}
 +ie^{-\{\frac12(y-y')+\frac12i(x'-x)\}}}{\surd\{(e^y+e^{-y})(e^{y'}+e^{-y'})\}}\\
 &=\frac{2\sin\{\tfrac12(x'-x)+\tfrac12i(y'-y)\}}
 {\surd\{(e^y+e^{-y})(e^{y'}+e^{-y'})\}}.
\end{aligned}
\]
\noindent [Ebenso gilt]
\[
\begin{aligned}
\cos\tfrac12D\,.\,e^{\frac12i(A-A')}
 &=\frac{\xi\xi+1}{2\xi}\cos\tfrac12(Y'-Y)
 +\frac{\xi\xi-1}{2\xi}\sin\tfrac12(Y'+Y)\\
 &=\frac{\xi\xi+1}{2\xi}\cos\tfrac12(p-p')
 +\frac{\xi\xi-1}{2\xi}\cos\tfrac12(p+p')\\
 &=\xi\cos\tfrac12p\cos\tfrac12p'
 +\xi^{-1}\sin\tfrac12p\sin\tfrac12p'\\
 &=\frac{\xi}{\surd\{(1+e^{-2y})(1+e^{-2y'})\}}
 +\frac{\xi^{-1}}{\surd\{(1+e^{2y})(1+e^{2y'})\}}\\
 &=\frac{\xi e^{\frac12(y+y')}+\xi^{-1}e^{-\frac12(y+y')}}
 {\surd\{(e^y+e^{-y})(e^{y'}+e^{-y'})\}}\\
 &=\frac{e^{\frac12(y+y')+\frac12i(x'-x)}
 +e^{-\{\frac12(y+y')+\frac12i(x'-x)\}}}{\surd\{(e^y+e^{-y})(e^{y'}+e^{-y'})\}}\\
 &=\frac{2\cos i\{\tfrac12(y+y')+\tfrac12i(x'-x)\}}
 {\surd\{(e^y+e^{-y})(e^{y'}+e^{-y'})\}}.
\end{aligned}
\]
\noindent [Man setze im einen Fall]
\[
\begin{aligned}
x'-x&=\delta\cos\alpha\\
y'-y&=\delta\sin\alpha,
\end{aligned}
\]
\noindent\hfill \(17^{*}\)\hspace*{.13\linewidth}\par
\endgroup
